% ===========================================================================
\section{Fundamenta\c{c}\~ao te\'orica}
\label{sec:fundamentacao}
% ===========================================================================

Esta se\c{c}\~ao estabelece os quatro corpos de conhecimento mobilizados pelo trabalho:
representa\c{c}\~ao e subdivis\~ao de superf\'icies, sombreamento fisicamente baseado,
gest\~ao de cor na cadeia de imagem, e agentes de intelig\^encia artificial que operam
ferramentas externas.

\subsection{Representa\c{c}\~ao de superf\'icies e malhas poligonais}

Uma superf\'icie de forma livre admite duas fam\'ilias principais de representa\c{c}\~ao. A
primeira \'e a param\'etrica cont\'inua, cujo caso mais geral \'e a NURBS, em que a superf\'icie
\'e definida por pontos de controle, pesos e vetores de n\'os \cite{piegl1997, farin2002}. A
segunda \'e a discreta: a superf\'icie \'e aproximada por uma malha poligonal de
conectividade expl\'icita \cite{mortenson1997, bbotsch2010}.

A escolha entre as duas fam\'ilias tem consequ\^encias pr\'aticas. A representa\c{c}\~ao
param\'etrica oferece continuidade anal\'itica controlada e \'e a base dos sistemas de
CAD/CAM \cite{shah1995}. A representa\c{c}\~ao discreta, por outro lado, \'e a que os
motores de render de tempo real e a maioria dos \emph{pipelines} de conte\'udo
tridimensional consomem, e sua qualidade passa a depender da \emph{topologia}: da
distribui\c{c}\~ao de faces, da val\^encia dos v\'ertices e da dire\c{c}\~ao das cadeias de
arestas \cite{bbotsch2010}.

Na pr\'atica automotiva, a malha de controle \'e constru\'ida com densidade baixa e
posteriormente suavizada por subdivis\~ao. \'E essa combina\c{c}\~ao que permite ao
modelador concentrar-se na \emph{inten\c{c}\~ao} de forma --- o perfil de uma se\c{c}\~ao
transversal, a posi\c{c}\~ao de uma aresta viva --- e delegar o refinamento ao algoritmo.

\subsection{Subdivis\~ao de superf\'icies e o papel da topologia}

O esquema de Catmull--Clark \cite{catmull1978} generaliza as B-splines c\'ubicas para
malhas de topologia arbitr\'aria: cada face \'e dividida por pontos de face e aresta, e cada
v\'ertice \'e reposicionado por uma m\'edia ponderada de seus vizinhos. O esquema de Loop
\cite{loop1987} faz o an\'alogo para malhas triangulares. Ambos convergem para uma
superf\'icie limite cuja continuidade depende da val\^encia dos v\'ertices; Stam
\cite{stam1998} demonstrou como avali\'a-la exatamente em par\^ametros arbitr\'arios, e
Halstead \emph{et al.} \cite{halstead1993} mostraram como interpolar atributos de forma
justa sobre a superf\'icie limite --- resultado central para a leitura de reflexos.

Tr\^es propriedades da topologia determinam o comportamento sob subdivis\~ao. A primeira \'e a
\emph{predomin\^ancia de quads}. Faces de quatro lados produzem, a cada n\'ivel, uma
subdivis\~ao regular; tri\^angulos e n-gons geram v\'ertices extraordin\'arios que puxam a
superf\'icie e criam ondula\c{c}\~oes nos reflexos \cite{bbotsch2010}. A segunda \'e a
\emph{densidade local}. Uma varia\c{c}\~ao de curvatura de raio pequeno --- uma aresta viva,
um vinco --- s\'o sobrevive \`a suaviza\c{c}\~ao se estiver sustentada por an\'eis de
arestas pr\'oximos, inseridos deliberadamente para ``segurar'' a superf\'icie limite. A
terceira \'e o \emph{fluxo de arestas}: cadeias alinhadas \`a dire\c{c}\~ao dominante da
forma distribuem a curvatura de modo previs\'ivel e facilitam a deforma\c{c}\~ao posterior.

Essas tr\^es propriedades t\^em uma consequ\^encia de projeto que este trabalho explora
diretamente. Aceitar uma opera\c{c}\~ao booleana para recortar um arco de roda introduz,
sobre a aresta vis\'ivel do para-lama, exatamente os elementos que degradam o reflexo:
n-gons, tri\^angulos e v\'ertices extraordin\'arios. Impor a restri\c{c}\~ao
\emph{no algoritmo} --- isto \'e, construir apenas estruturas que n\~ao os produzam ---
\'e uma postura diferente de corrigir a malha depois.

\subsection{Sombreamento fisicamente baseado}

A equa\c{c}\~ao de render de Kajiya \cite{kajiya1986} formaliza a forma\c{c}\~ao de imagem
como uma integral de transporte de radi\^ancia. Modelos de reflex\^ao que conservam energia
--- de Phong \cite{phong1975} e Blinn \cite{blinn1977}, passando por Cook--Torrance
\cite{cook1982}, Schlick \cite{schlick1994}, modelos anisotr\'opicos \cite{ward1992,
ashikhmin2000} e microfacetas para refra\c{c}\~ao \cite{walter2007} --- s\~ao
aproxima\c{c}\~oes dessa integral, e sua escolha determina a plausibilidade f\'isica do
resultado \cite{pharr2016}.

A pr\'atica moderna convergiu para um modelo \'unico parametrizado, no qual cor base,
metalicidade, rugosidade, refletividade, transmiss\~ao e camadas adicionais de verniz s\~ao
par\^ametros independentes \cite{burley2012, karis2013}. Para pintura automotiva, a
camada relevante \'e o \emph{clear coat}: uma superf\'icie diel\'etrica de rugosidade muito
baixa sobre a base pigmentada, que produz o especular n\'itido e o afastamento de
tonalidade caracter\'isticos. Como a refletividade de um diel\'etrico cresce com o
\^angulo de incid\^encia, o escurecimento das bordas em \^angulos rasantes emerge do
pr\'oprio modelo f\'isico, e n\~ao de um ajuste arbitr\'ario.

\subsection{Cor e cadeia de imagem}

O transporte de radi\^ancia \'e linear na energia transportada, enquanto a percep\c{c}\~ao
de cor e a codifica\c{c}\~ao de imagens n\~ao s\~ao \cite{fairchild2013}. Duas
consequ\^encias pr\'aticas decorrem disso. A primeira \'e que cores definidas em
especifica\c{c}\~oes visuais (sRGB) precisam ser convertidas para o espa\c{c}\~o linear
antes de alimentarem o sombreamento; a omiss\~ao dessa convers\~ao produz um erro sutil, em
que a imagem \emph{parece} plaus\'ivel mas tem satura\c{c}\~ao sistematicamente
deslocada. A segunda \'e que o resultado precisa ser mapeado de volta para um espa\c{c}\~o
de exibi\c{c}\~ao por uma curva de mapeamento de tons \cite{reinhard2002}; a escolha dessa
curva comp\~oe a imagem final tanto quanto o material. Comparar o tom de um render com o
tom de uma fotografia sem controlar as duas cadeias \'e, em rigor, uma
compara\c{c}\~ao inv\'alida.

\subsection{Agentes e protocolos de contexto}

Modelos de linguagem de grande porte, baseados em mecanismos de aten\c{c}\~ao
\cite{vaswani2017} e treinados em escala \cite{brown2020}, demonstraram capacidade de
racioc\'inio em etapas \cite{wei2022} e de uso de ferramentas externas
\cite{schick2023, patil2023, qin2023}. A integra\c{c}\~ao entre racioc\'inio e
a\c{c}\~ao \cite{yao2023}, a recupera\c{c}\~ao de contexto \cite{lewis2020} e a
coordena\c{c}\~ao entre agentes \cite{wu2023} completam o quadro. Essas contribui\c{c}\~oes
resolviam, por\'em, a integra\c{c}\~ao no n\'ivel da \emph{aplica\c{c}\~ao}.

O MCP \cite{mcpspec2024} acrescenta uma camada de protocolo: um servidor exp\~oe
ferramentas de forma descobr\'ivel e o cliente --- um agente --- as invoca por mensagens
JSON-RPC~2.0 \cite{jsonrpc2020}, sem depend\^encia do fornecedor do modelo. Do ponto de
vista da engenharia, a propriedade relevante n\~ao \'e a capacidade de escrever c\'odigo,
que existe em qualquer fluxo por \emph{script}, e sim a \emph{serializa\c{c}\~ao da
inten\c{c}\~ao}: cada chamada fica registrada como uma mensagem com nome de ferramenta e
argumentos, o que converte a sess\~ao de modelagem em um registro inspecion\'avel.
